\section{VOO Syntax Sugar vs.\ Raw Vanilla Tcl}
\label{app:syntax}

This appendix demonstrates the value of VOO's syntax sugar by comparing a \texttt{Point}
class written with VOO's declarative syntax against the equivalent raw Tcl that VOO
generates internally.

\subsection*{C.1 Class Declaration and Field Indices}

\begin{lstlisting}[language=Tcl, caption={VOO syntax sugar: class and fields}]
voo::class Point {
    public {
        double_t x 0.0
        double_t y 0.0
        string_t name "point"
    }
}
\end{lstlisting}

\begin{lstlisting}[language=Tcl, caption={Raw Vanilla Tcl generated by VOO}]
namespace eval Point {
    variable x 0
    variable y 1
    variable name 2

    variable __defaultObj [list 0.0 0.0 "point"]
    variable __fields [list x y name]
}
\end{lstlisting}

\subsection*{C.2 Constructors}

\subsubsection*{C.2.1 Positional Constructor}
\begin{lstlisting}[language=Tcl]
# VOO: auto-generated Point::new
set p [Point::new 1.5 2.5 "A"]

# Raw (generated):
proc Point::new {x y name} { return [list $x $y $name] }
\end{lstlisting}

\subsubsection*{C.2.2 No-Argument Constructor}
\begin{lstlisting}[language=Tcl]
# VOO:
set p [Point::new()]

# Raw (generated):
proc Point::new() {} {
    variable __defaultObj
    return $__defaultObj
}
\end{lstlisting}

\subsubsection*{C.2.3 Named-Argument Constructor}
\begin{lstlisting}[language=Tcl]
# VOO:
set p [Point::new.args -x 1.5 -y 2.5]

# Raw (generated):
proc Point::new.args {args} {
    variable __defaultObj
    set obj $__defaultObj
    if {[catch {dict size $args}]} {
        error "Constructor argument must be a list of '-<field> <value>' pairs"
    }
    dict for {key value} $args {
        if {[string index $key 0] ne "-"} {
            error "Constructor argument keys must start with '-', got '$key'"
        }
        set field [string range $key 1 end]
        set setter set.$field
        if {[info commands $setter] ne ""} {
            $setter obj $value
        } else {
            set setter my.set.$field
            if {[info commands $setter] ne ""} {
                $setter obj $value
            } else {
                error "Unknown field option: $field"
            }
        }
    }
    return $obj
}
\end{lstlisting}

\subsection*{C.3 Getters}
\begin{lstlisting}[language=Tcl]
# VOO: auto-generated
set x [Point::get.x $p]

# Raw (generated per field):
proc Point::get.x {this} { variable x; return [lindex $this $x] }
proc Point::get.y {this} { variable y; return [lindex $this $y] }
proc Point::get.name {this} { variable name; return [lindex $this $name] }
\end{lstlisting}

\subsection*{C.4 Setters}
\begin{lstlisting}[language=Tcl]
# VOO: auto-generated
Point::set.x p 3.14

# Raw (generated per field):
proc Point::set.x {thisVar value} {
    variable x
    upvar $thisVar this
    lset this $x $value
}
\end{lstlisting}

Setters use \texttt{upvar} to receive the variable name, enabling in-place modification
with copy-on-write safety.

\subsection*{C.5 Updaters (Copy-on-Write Optimization)}
\begin{lstlisting}[language=Tcl]
# VOO: auto-generated
Point::update.x p temp { set temp [expr {$temp * 2}] }

# Raw (generated per field):
proc Point::update.x {thisVar tempVar body} {
    variable x
    upvar $thisVar this
    upvar $tempVar temp
    try {
        set temp [lindex $this $x]
        lset this $x {}
        uplevel $body
    } finally {
        lset this $x $temp
    }
}
\end{lstlisting}

\subsection*{C.6 Instance Methods}
\begin{lstlisting}[language=Tcl]
# VOO:
voo::class Point {
    method distance {} {
        set dx [get.x $this]; set dy [get.y $this]
        return [expr {sqrt($dx*$dx + $dy*$dy)}]
    }
}

# Raw (generated):
proc Point::distance {this} {
    variable x; variable y
    set dx [lindex $this $x]; set dy [lindex $this $y]
    return [expr {sqrt($dx*$dx + $dy*$dy)}]
}
\end{lstlisting}

\subsection*{C.7 Static Fields}
\begin{lstlisting}[language=Tcl]
# VOO:
voo::class Point {
    public { int_t -static count 0 }
}
set n [Point::class.get.count]

# Raw (generated):
namespace eval Point {
    variable count 0
    proc class.get.count {} { variable count; return $count }
    proc class.set.count {value} { variable count; set count $value }
}
\end{lstlisting}

\subsection*{C.8 Virtual Classes and Methods}

\subsubsection*{C.8.1 Virtual Base Class Declaration}
\begin{lstlisting}[language=Tcl, caption={VOO virtual base}]
voo::class Shape -virtual {
    public { double_t radius 1.0 }
    method area -virtual {} { return 0.0 }
}
\end{lstlisting}

\begin{lstlisting}[language=Tcl, caption={Raw Vanilla Tcl generated for virtual base}]
namespace eval Shape {
    # Index 0 is permanently reserved for the class namespace tag.
    variable radius 1          ;# field index = 1
    variable __defaultObj [list ::Shape 1.0]
    variable __fields [list radius]
    variable __voo_is_virtual_class 1

    proc new {radius} { return [list ::Shape $radius] }

    # base.area holds the original body for direct parent calls
    proc base.area {this} { return 0.0 }

    # area is a dispatcher
    proc area {this} {
        set __voo_cls [lindex $this 0]
        if {$__voo_cls ne [namespace current] && \
            [info commands ${__voo_cls}::area] ne {}} {
            return [${__voo_cls}::area $this]
        }
        return [base.area $this]
    }
}
\end{lstlisting}

\subsubsection*{C.8.2 Virtual Child Class Declaration}
\begin{lstlisting}[language=Tcl, caption={Child class overriding a virtual method}]
voo::class Circle -extends Shape {
    method area -override {} {
        return [expr {3.14159 * [get.radius $this] ** 2}]
    }
}
\end{lstlisting}

\begin{lstlisting}[language=Tcl, caption={Raw Vanilla Tcl generated for child class}]
namespace eval Circle {
    variable radius 1          ;# same index as Shape::radius
    variable __defaultObj [list ::Circle 1.0]
    variable __voo_is_virtual_class 1

    proc new {radius} { return [list ::Circle $radius] }

    proc base.area {this} {
        return [expr {3.14159 * [get.radius $this] ** 2}]
    }
    proc area {this} {
        set __voo_cls [lindex $this 0]
        if {$__voo_cls ne [namespace current] && \
            [info commands ${__voo_cls}::area] ne {}} {
            return [${__voo_cls}::area $this]
        }
        return [base.area $this]
    }
}
\end{lstlisting}

\subsubsection*{C.8.3 Calling the Parent Body from an Override}
\begin{lstlisting}[language=Tcl]
voo::class ColoredCircle -extends Circle {
    public { string_t color "red" }
    method area -override {} {
        # Add 10% for visual padding
        set base [Circle::base.area $this]
        return [expr {$base * 1.1}]
    }
}
\end{lstlisting}
